\chapter{复变函数的概念}

\problem{4.1}{写出下列复数的实部、虚部、模和幅角。
\begin{enumerate}
  \item $1-\mathrm{i}\sqrt{3}$ \qquad
  \item $\ln\mathrm{i}$ \qquad
  \item $(1+\mathrm{i})^{1-\mathrm{i}}$ \qquad
  \item $\sqrt[\mathrm{i}]{\mathrm{i}}$
  \item $\sin(a+\mathrm{i}b),\quad a,b\in R$ \qquad
  \item $\sqrt{a+\mathrm{i}b},\quad a,b\in R$
\end{enumerate}}
\answerarea

\problem{4.2}{写出下列复变函数 $W=f(z)=u(x,y)+\mathrm{i}v(x,y)$ 的实部 $u(x,y)$、虚部 $v(x,y)$、模 $|f(z)|$ 和幅角 $\arg f(z)$。
\begin{enumerate}
  \item $\mathrm{e}^z$ \qquad
  \item $\mathrm{e}^{\mathrm{i}z}$ \qquad
  \item $\mathrm{e}^{\mathrm{i}\sin x},\quad x\in R$
  \item $z^3$ \qquad
  \item $\sqrt{\mathrm{i}+\mathrm{i}2x\sqrt{x^2-1}},\quad x\in R$
\end{enumerate}}
\answerarea
\clearpage

\problem{4.3}{求下列方程的解，并把解在复平面上用图形表示出来，说明它们的几何意义。
\begin{enumerate}
  \item $|z|=1$ \qquad
  \item $|z|<1$
  \item $|z|\geqslant 1$ \qquad
  \item $1<|z|<2$
  \item $|z+\mathrm{i}|\leqslant |z-\mathrm{i}|$ \qquad
  \item $|z-a|=|z-b|\quad a,b\text{ 为复平面上的点}$
  \item $0<\arg(z-1)<\dfrac{\pi}{4}$ \qquad
  \item $\operatorname{Re}\!\left(\dfrac{1}{z}\right)=\dfrac{1}{2}$
  \item $|z_1+z_2|^2+|z_1-z_2|^2=2|z_1|^2+2|z_2|^2$
  \item $\sin z=0$ \qquad
  \item $\cos z=0$
  \item $\sh z=0$ \qquad
  \item $\ch z=0$
\end{enumerate}}
\answerarea

\problem{4.4}{试证明：
\begin{enumerate}
  \item $|\sin z|\geqslant \sin x$ \qquad
  \item $|\cos z|\geqslant |\cos x|$
\end{enumerate}}
\answerarea
\clearpage

\problem{4.5}{对于在实变量 $x$ 中成立的关系式
\[
  \sin(x_1+x_2) = \sin x_1\cos x_2 + \sin x_2\cos x_1
\]
请证明这一关系式在复变量的情况下也成立，即
\[
  \sin(z_1+z_2) = \sin z_1\cos z_2 + \sin z_2\cos z_1
\]}
\answerarea

\problem{4.6}{对下列复变函数作出 Riemann 面的示意图，并指出函数的枝点及其阶数，同时说明你所作的 Riemann 面的割线的取法。
\begin{enumerate}
  \item $\sqrt{(z-a)(z-b)}$ \qquad
  \item $\sqrt[3]{(z-a)(z-b)}$ \qquad
  \item $\ln(z-a)$
  \item $\sqrt{(z-a)(z-b)(z-c)}$
\end{enumerate}
其中 $a,b,c$ 为互不相等的复常数。}
\answerarea
\clearpage
